\section*{Notation}
\addcontentsline{toc}{section}{Notation}

\emph{Typesetter's note:} the source conventions below are presented with the
modernizations specified in \texttt{NOTATION.md}. In particular, spacetime
components are ordered $(0,1,2,3)$ while retaining Weinberg's mostly-plus
metric, states use Dirac notation, and Batalin--Vilkovisky antifields use
\(\ddagger\).

Latin indices \(i,j,k,\ldots\) generally run over the three spatial coordinate
labels \(1,2,3\). Where specifically indicated, they may run over the values
\(1,2,3,4\), with the source convention \(x^4=it\).

Greek indices \(\mu,\nu,\ldots\) from the middle of the Greek alphabet
generally run over the four spacetime coordinate labels \(0,1,2,3\), with
\(x^0\) the time coordinate.

Lowercase Latin indices \(a,b,c,\ldots\) from the beginning of the alphabet
generally run over the generators of a symmetry algebra, including
gauge-adjoint indices. If these letters already label another independent
family in the same derivation, a later lowercase Latin run such as
\(r,s,t,\ldots\) is used to avoid a collision. Early Greek letters remain
available for state, spinor, coupling, and other non-adjoint labels.

When a quantity carries both a gauge-adjoint index and spacetime indices,
their opposing scripts are set together, as in
\(F^a_{\mu\nu}\) and \(A_a^\mu\);
an empty group is not inserted between the scripts.

Repeated indices are generally summed, unless otherwise indicated.

The spacetime metric is mostly plus,
\[
  \eta_{\mu\nu}=\operatorname{diag}(-1,+1,+1,+1).
\]

The d'Alembertian is defined as
\[
  \Box=\eta^{\mu\nu}
  \frac{\partial^2}{\partial x^\mu\partial x^\nu}
  =\nabla^2-\frac{\partial^2}{\partial t^2},
\]
where \(\nabla^2\) is the spatial Laplacian
\(\partial^2/\partial x^i\partial x^i\).

The Levi--Civita tensor \(\epsilon^{\mu\nu\rho\sigma}\) is defined as the
totally antisymmetric quantity with \(\epsilon^{0123}=+1\).

Spatial three-vectors and internal three-vectors are indicated by letters in
boldface, using \(\mathbf{x}\) for Latin symbols and \(\boldsymbol{\pi}\) for
Greek symbols. Arrow accents are not used for vectors.

A hat over any vector indicates the corresponding unit vector; thus
\(\widehat{\mathbf V}=\mathbf V/\lvert\mathbf V\rvert\).

A dot over any quantity denotes the time derivative of that quantity.

Dirac matrices \(\gamma^\mu\) are defined so that
\[
  \gamma^\mu\gamma^\nu+\gamma^\nu\gamma^\mu=2\eta^{\mu\nu}.
\]
Also,
\[
  \gamma_5=i\gamma^0\gamma^1\gamma^2\gamma^3,
  \qquad
  \beta=i\gamma^0.
\]

The step function \(\theta(s)\) has the value \(+1\) for \(s>0\) and \(0\) for
\(s<0\).

The complex conjugate, transpose, and Hermitian adjoint of a matrix or vector
\(A\) are denoted \(A^*\), \(A^{\mathsf T}\), and \(A^\dagger\), respectively.
The Hermitian adjoint of an operator \(O\) is denoted \(O^\dagger\), except
where an asterisk is used to emphasize that a vector or matrix of operators is
not transposed. \(+\mathrm{H.c.}\) or \(+\mathrm{c.c.}\) at the end of an
expression indicates the addition of the Hermitian adjoint or complex
conjugate of the foregoing terms. A bar on a Dirac spinor \(u\) is defined by
\(\bar u=u^\dagger\beta\). The antifield of a field \(x^n\) in the
Batalin--Vilkovisky formalism is denoted \(x_n^{\ddagger}\), rather than
\(x_n^*\), to distinguish it from the ordinary complex conjugate or the
antiparticle field.

States use Dirac notation. In particular, asymptotic states are written
\(\InKet{\alpha}\) and \(\OutKet{\beta}\), with the
labels \(\mathrm{in}\) and \(\mathrm{out}\) outside the ket or bra. Editorial
exercise material uses the helper macros documented in \texttt{NOTATION.md}
to enforce this convention.

The positive energy of a momentum mode is denoted
\(\omega_{\mathbf k}\), or \(\omega_k\) when the momentum label is not
bold. The letter \(E\) is reserved for total energies, energy eigenvalues, or
quantities conventionally named \(E\).

Units are usually used with \(\hbar\) and the speed of light taken to be unity.
Throughout, \(-e\) is the rationalized charge of the electron, so that the
fine-structure constant is
\[
  \alpha=\frac{e^2}{4\pi}\simeq\frac{1}{137}.
\]

Numbers in parentheses at the end of quoted numerical data give the
uncertainty in the last digits of the quoted figure. Where not otherwise
indicated, experimental data are taken from `Review of Particle Properties,'
\emph{Phys.\ Rev.} \textbf{D50}, 1173 (1994).
